% s01-intro.tex —— 导引（原文 Introduction, p.1–3）
% 原文：Laurent Fargues, Geometrization of the local Langlands correspondence: an overview, 2016

\section*{导引}
\addcontentsline{toc}{section}{导引}
\markright{导引}
\srcpage{1}

本文是一份对某个猜想的非正式概述；该猜想把 p-进 Hodge 理论、几何 Langlands 纲领与局部 Langlands 对应联系起来。
更精确、也更技术性的文章将在稍后发表（[7]）。

固定局部域 $E$ 上的拟分裂约化群 $G$，其中 $E=\FF_q((\pi))$ 或 $[E:\QQ_p]<+\infty$，$\FF_q=\oo_E/\pi$。
固定离散 Langlands 参数 $\varphi:W_E\to{}^{L}G$，这里 ${}^{L}G$ 是 $G$ 的 $\ell$-进 Langlands 对偶群，$\ell\neq p$。
猜想对每个这样的 $\varphi$ 配一个 $S_\varphi$-等变的 Hecke 特征层 $\cF_\varphi$，
它住在我们与 Fontaine 合作构造并研究过的曲线上的 $G$-丛叠里（[10]）。\yizhu{%
Hecke 特征层：沿 Hecke 对应具有本征性的 $\ell$-进层；本征值取局部系统 $r_\mu\circ\varphi$。%
}{A：原文导引 p.1；I：本库 SY02-19、SY02-21（「Hecke 特征层」）}

更确切地说，给定 $\FF_q$-完满空间 $S$，可以构造一个 $E$-adic 空间 $X_S$。
应当把它想成曲线族 $(X_{k(s)})_{s\in S}$，其中 $X_{k(s)}$ 是完满域 $k(s)$ 对应的 p-进 Hodge 理论基本曲线的 adic 版本（[10]，[9]）。
当 $S=\Spa(R,R^+)$ 是 affinoid 完满空间时，它有一个概形式的对应物：一个 $E$-概形，
并带有该 $E$-概形上的 $G$-丛与 $X_S$ 上的 $G$-丛之间的 GAGA 等价。\yizhu{%
完满空间（perfectoid space）是 Scholze 引入的一类 adic 空间，局部由“完满”的 affinoid 代数 $\overline{\QQ}_p\langle T_1^{\pm1/p^\infty},\dots\rangle$ 图卡给出；本文的完满空间不固定在一个完满基域上。%
}{A：原文 §1.1 p.3；B：nLab "perfectoid space"；I：本库 SY02-16}
$G$-丛叠就是叠
\[
\Bun_G:\ S\longmapsto X_S\ \text{上}\ G\text{-丛的群胚}.
\]

\srcpage{2}

这是 pro-étale 拓扑下的叠。在 §2.1 中我们将说明，如何通过在局部“光滑”拓扑下的一套完满图卡，在这个叠上放一个几何结构。
光滑态射这个概念目前仍然模糊，但我们试图说服读者（以及我们自己）：Quot 格式的技巧可以搬到我们的情形里。
从这个观点看，代数空间的类比物就是 Scholze 引入的菱形概念（[27]）。\yizhu{%
菱形（diamond）：$\mathrm{Perf}_{\FF_p}$ 中按 pro-étale 拓扑取的代数空间，即完满空间按 pro-étale 等价关系取商所得的几何对象；正常刚性空间都嵌入菱形范畴。%
}{A：原文 §1.3.1 p.8；B：nLab "diamond (pro-étale)"; I：本库 SY02-26 裁定}
人们希望能在这个几何语境里有好的反常 $\ell$-进层理论。
Caraiani–Scholze 近来的工作（[2]）表明这样的反常层理论可能存在，参见本文 §7 关于局部/整体相容的部分。
还有另一个观点可以在 $\Bun_G$ 上放几何结构，即通过 Beauville–Laszlo 单值化，参见 §3.2。

让我们记下
\begin{itemize}
\item 当 $E=\FF_q((\pi))$ 时，$\Spa(E)^\diamond=\Spa(E)$；
\item 当 $E\mid\QQ_p$ 时，$\Spa(E)^\diamond$ 是 Scholze 定义的 untilt 层（[27]）。
\end{itemize}
任何态射 $S\to\Spa(E)^\diamond$ 都定义一个 Cartier 除子
\[
S^\sharp\hookrightarrow X_S,
\]
其中当 $E=\FF_q((\pi))$ 时 $S^\sharp=S$；当 $E\mid\QQ_p$ 时 $S^\sharp$ 是 $S$ 相应的 untilt。\yizhu{%
untilt：完满空间 $S$ 的一个 untilt 是一对 $(S^\sharp,\iota)$，其中 $S^\sharp$ 是 $E$ 上的完满空间、$\iota:S\xrightarrow{\sim}S^{\sharp,\flat}$ 是到其倾斜的同构。中文译名未定，保留原词。%
}{A：原文 §1.3.1 定义 1.8 p.8；I：本库 SY02-26 裁定（保留原词）}
当 $S=\Spa(R,R^+)$ 时，$X_S$ 沿这个 Cartier 除子的形式完备化就是 Fontaine 环 $B_{\mathrm{dR}}^+(R^\sharp)$ 的形式谱。
这一切让我们可以定义 Hecke 对应
\begin{equation}\label{eq:0.1}
\begin{tikzcd}
& \Hck^{\le\mu} \arrow[dl] \arrow[dr] & \\
\Bun_G & & \Bun_G\times\Spa(E)^\diamond
\end{tikzcd}
\end{equation}
其中右向态射是 $B_{\mathrm{dR}}$-仿射 Grassmann 簇中与 $\mu$ 相联系的闭 Schubert 胞腔里的局部平凡纤维化（[27]）。
当 $E=\FF_q((\pi))$ 时，这个 $B_{\mathrm{dR}}$-仿射 Grassmann 簇就是 Pappas–Rapoport 扭仿射 Grassmann 簇的刚性解析对应物（[21]）。

猜想说：$\cF_\varphi$ 应当是这些 Hecke 对应的特征层，特征值等于局部系统 $r_\mu\circ\varphi$，
其中 $r_\mu$ 是 ${}^{L}G$ 的与 $\mu$ 对偶的通常表示。

若 $B(G)$ 是 $G(\breve E^{\mathrm{un}})$ 中 $\sigma$-共轭类的 Kottwitz 集（[18]），则 [6] 的主要定理给出一个双射\yizhu{%
$G(\breve E^{\mathrm{un}})$ 中两个元素 $\sigma$-共轭，指存在 $g$ 使 $b'=g b\sigma(g)^{-1}$；Kottwitz 集 $B(G)$ 就是这些类的集合，它同时分类曲线上的 $G$-丛。%
}{A：原文导引 p.2；B：nLab "Kottwitz set"；I：本库 SY02-16}
\[
B(G)\xrightarrow{\ \sim\ }|\Bun_{G,\FF_q}|.
\]
具体地说：设 $\cE$ 是 $X_S$ 上的 $G$-丛，$S$ 是 $\FF_q$ 上的完满空间。
对每个几何点 $\bar s\to S$，$\cE|_{X_{k(\bar s)}}$ 由 Kottwitz 集 $B(G)$ 的一个元素分类。
这就定义了映射 $|S|\to B(G)$，它记录 $\cE$ 在每个几何纤维上的限制。

若 $[b]\in B(G)$ 是基本的，则 $J_b$（$b$ 的 $\sigma$-中心化子）是 $G$ 的内形式。
并且 $\Bun_{G,\FF_q}$ 对应点处的剩余 gerbe 同构于 pro-étale $J_b(E)$-torsor 的分类空间。
于是 $\cF_\varphi$ 在该点的茎给出 $J_b(E)$ 的一个光滑 $\overline{\QQ}_\ell$-表示，
连同 $S_\varphi$ 的一个交换作用；后者来自 $S_\varphi$ 在 $\cF_\varphi$ 上的作用。
猜想说：$S_\varphi$ 在茎上的这个线性作用，恰好切出 $J_b(E)$ 的一个 $L$-包，
对应的是与 $G$ 的扩充纯内形式 $J_b$ 相联系的局部 Langlands 对应。

这样的扩充纯内形式出现在 Kaletha 的工作里（例如 [13]，[14]）。
它们通过嵌入 $H^1(E,G)\subset B(G)$ 推广了 Vogan 的纯内形式。
这个猜想的动机之一在于：拉回到曲线上之后，这些约化群概形 $J_b$ 变成 $G$ 的纯内形式，于是我们又落回 Vogan 的语境。

还要提到 $\cF_\varphi$ 的另一个性质，即所谓的特征层性质。
事实是：$\cF_\varphi$ 不应当是 $\Bun_G$ 上的反常层，而应当是 $\Bun_G\otimes\FF_q$ 上的 Weil-反常层。

有一个映射 $\{G(E)\}_{\mathrm{ell}}\to B(G)_{\mathrm{basic}}$。
若 $\delta\in G(E)$ 是椭圆元素，则 $\Bun_G$ 中对应的点 $x_\delta$ 定义在 $\FF_q$ 上，
因此前一个茎 $x_\delta^*\cF_\varphi$ 应当带上一个 Frobenius 作用。\srcpageinl{3}特征层性质蕴涵：函数
\begin{equation}\label{eq:0.2}
\begin{aligned}
\{G(E)\}_{\mathrm{ell-reg}} &\longrightarrow \overline{\QQ}_\ell,\\
\{\delta\} &\longmapsto \operatorname{Tr}\big(\mathrm{Frob};\,x_\delta^*\cF_\varphi\big)
\end{aligned}
\end{equation}
就是与 $\varphi$ 相联系的 $G(E)$ 稳定分布在椭圆正则子集上的限制。\yizhu{%
稳定分布：在 $G(E)$ 的轨道积分意义下不变的分布；特征层性质说，$\cF_\varphi$ 的 Frobenius 迹在椭圆正则元上给出该稳定分布。%
}{A：原文导引 p.3；B：nLab "stable distribution (Langlands)"}

本文的结构如下：
\begin{itemize}
\item §1、§2、§3 解释猜想中出现的对象的构造与几何：曲线 $X_S$、叠 $\Bun_G$、相应的 Hecke 叠以及 $B_{\mathrm{dR}}$-仿射 Grassmann 簇。
  其间我们会快速回忆 Scholze 的一些构造（[27]），例如他的菱形理论。
\item §4 给出猜想的精确陈述。
\item §5 的目的是更详细地解释猜想中出现的特征层性质。
\item §6 是 §7 与 §8 的预备。我们说明：不等特征下 p-可除群的 $F$-等晶体的 Hodge 滤过，定义了从 p-可除群预叠到 Hecke 叠的预叠态射。
  从这个观点看，把 Hecke 叠定义为对应的两个态射 $\overleftarrow{h}$ 与 $\overrightarrow{h}$，可以分别解释为 de Rham 周期态射与 Hodge–Tate 周期态射。
\item §7 解释猜想的局部/整体相容部分。它说：我们猜想的局部反常层 $\cF_\varphi$ 与 Caraiani–Scholze 的反常层 $R\pi_{\mathrm{HT},*}\overline{\QQ}_\ell$ 相容，
  后者由 Hodge 型 Shimura 簇在其 Hodge–Tate 周期空间上的相对上同调构造。
\item §8 解释为什么我们的猜想蕴涵 Kottwitz 猜想，后者描述 Rapoport–Zink 定义的 p-可除群模空间上同调的离散部分（[23]）。
  我们先解释 Lubin–Tate/Drinfeld 情形（§8.1）：对这个关键情形的研​究，正是引导作者得到其猜想的主要因素。
\item §9 对环面验证猜想。此时猜想由局部类域论推出。
\end{itemize}

\begin{flushleft}
\itshape 致谢：作者感谢 Kiran Kedlaya、Peter Scholze，以及 2014 年秋在 MSRI 举行的项目“New Geometric Methods in Number Theory and Automorphic Forms”的所有参与者。
\end{flushleft}
